%=============================================================================!
% 16.11  FREE ENERGY FROM A HISTOGRAM                                         !
%=============================================================================!
\section{Free energy from a histogram}
\label{sec:ob-fes}

A trajectory spends more time in some arrangements than in others. At
equilibrium these populations measure free-energy differences. A
histogram of a distance, angle or position therefore gives a free-energy
profile along that coordinate, provided the run has sampled it adequately.
We test this construction on a surface and a four-bead chain, where the
answer is known, then examine why rare crossings make the uncertainty
hard to estimate.

\subsection*{Free energy along a coordinate}

In the canonical ensemble the positions have the density $\mathcal{P}
\propto \mathrm{e}^{-\beta U}$ (\cref{sec:sm-canonical}). For a coordinate
$x$ computed from the positions, the density of its values is the total
weight of all arrangements with that value, and writing it as a Boltzmann
factor defines the free energy along $x$:
\begin{equation}
   \mathcal{P}(x) \propto \mathrm{e}^{-\beta F(x)} , \qquad
   F(x) = -\kB T\ln\mathcal{P}(x) + \text{constant} ,
   \label{eq:ob-fes}
\end{equation}
the free energy $-\kB T\ln Z$ of \cref{sec:sm-canonical} for the
arrangements held at that $x$. If $x$ is every coordinate of the system,
$F = U$ up to a constant. If it is fewer, $F$ includes the number of
arrangements that share each value, an entropy, and can differ from any
energy. A free-energy profile from a run is a histogram of $x$, as for
$g(r)$, turned into $F$ by \cref{eq:ob-fes}. Its error comes from blocks,
as in Chapter~15: the histogram of each block, normalised to a density;
the mean and the standard error of each bin over the blocks; and, since a
small change $\delta\mathcal{P}$ changes $\ln\mathcal{P}$ by
$\delta\mathcal{P}/\mathcal{P}$, an error of $\kB T\,\delta\mathcal{P}/
\mathcal{P}$ in $F$. \texttt{mdlab.analysis.landscape} does this for one
coordinate or several:

\begin{code}
def free_energy(
    values: ArrayLike,
    bins: int | ArrayLike | list,
    temperature: float,
    n_blocks: int = 10,
    value_range: list | None = None,
) -> tuple[list[NDArray], NDArray, NDArray]:
    x = np.asarray(values, dtype=float)
    x = x[:, None] if x.ndim == 1 else x
    edges = np.histogramdd(x, bins=bins, range=value_range)[1]
    blocks = np.array_split(x, n_blocks)
    density = np.array(
        [np.histogramdd(b, bins=edges, density=True)[0] for b in blocks]
    )
    p = density.mean(axis=0)
    dp = density.std(axis=0, ddof=1) / np.sqrt(n_blocks)
    kt = KB * temperature
    with np.errstate(divide="ignore", invalid="ignore"):
        f = -kt * np.log(p)
        df = kt * dp / p
    f -= np.min(f[np.isfinite(f)])
    middles = [0.5 * (e[1:] + e[:-1]) for e in edges]
    return middles, f, df
\end{code}

\subsection*{An atom on a surface}

A lithium atom on the model surface of \cref{sec:en-surface} has two
coordinates, and its position is all of them, so $F(x, y)$ must equal
$U(x, y)$. 1000 atoms at \SI{600}{\kelvin} under Langevin friction of
\SI{1}{\per\pico\second}, every \SI{50}{\femto\second} for
\SI{200}{\pico\second}, folded into one cell of the hollows, give the map
of \cref{fig:ob-fes}(a). A profile is tested against an exact one by the
mean of its squared misses divided by their squared errors. Each such
ratio has the mean 1 when the errors are right and, for misses of
Gaussian density, the variance $\ens{z^4} - 1 = 2$, so over $n$ independent bins the
mean scatters about 1 by $\sqrt{2/n}$; a mean of $c$ says that the
errors understate the misses by about $\sqrt c$. Below
\SI{0.3}{\electronvolt}, over 1400 bins, $F$ less $U$ gives 1.20, where
independent bins with correct errors would give $1.00 \pm 0.04$.
Histogram bins are correlated, so this spread is only a guide: it is constant to within about
1.1 times its errors, and the largest miss is
\SI{15.2}{\milli\electronvolt}.

The distance $r_{\mathrm h}$ from the nearest hollow is one coordinate of the two,
and $F(r_{\mathrm h})$ is not $U$. The positions at a distance $r_{\mathrm h}$ fill a circle of
circumference $2\pi r_{\mathrm h}$, so the density of $r_{\mathrm h}$ is
\begin{equation*}
   \mathcal{P}(r_{\mathrm h}) \propto r_{\mathrm h}\oint\mathrm{e}^{-\beta U(r_{\mathrm h}, \theta)}\,\dd\theta ,
   \qquad
   F(r_{\mathrm h}) = -\kB T\ln\Bigl[r_{\mathrm h}\oint\mathrm{e}^{-\beta U(r_{\mathrm h}, \theta)}\,\dd\theta
   \Bigr] ,
\end{equation*}
with $\theta$ the angle around the hollow, valid out to $a/2 =
\SI{1.23}{\angstrom}$, where the circle still lies in the region nearer
that hollow than any other. The factor $r_{\mathrm h}$ is the number of positions at
each distance, and it makes $F$ rise as $r_{\mathrm h} \to 0$, where $U$ is lowest:
there are few places that close to the centre, as there are few
velocities with all three components near zero in Maxwell's density of
speeds (\cref{sec:sm-maxwell}). The run's $F(r_{\mathrm h})$ follows
this within its errors, its squared misses over squared errors averaging
0.60 over 40 bins (\cref{fig:ob-fes}b), and its least value lies at $r_{\mathrm h} =
\SI{0.225}{\angstrom}$ from the hollow. For small $r_{\mathrm h}$, $U \approx
\frac12 kr_{\mathrm h}^2$ with $k = 3U_{\mathrm b}\lvert\vb{g}\rvert^2/8 =
2\pi^2U_{\mathrm b}/a^2$ (\cref{sec:ob-vdos}), so $\mathcal{P}
\propto r_{\mathrm h}\,\mathrm{e}^{-\beta kr_{\mathrm h}^2/2}$, whose largest value is at $r_{\mathrm h} =
\sqrt{\kB T/k} = \SI{0.230}{\angstrom}$. At the smallest $r_{\mathrm h}$, $F$ is
\SI{116}{\milli\electronvolt} above its least. A free-energy profile along
a coordinate carries the count of arrangements at each value, and a
minimum of $F$ need not be a minimum of $U$.

\subsection*{A dihedral angle}

The free energy of a dihedral angle is how the TrajCast paper judges the
slow motions of a molecule\cite{Thiemann2026}, and a model chain has an
exact answer for it. The chain has four beads of \SI{14.5}{\amu}, each
standing for a CH$_2$ or CH$_3$ group, joined by springs in the force
field of \cref{eq:pe-opls}: bonds with $\kappa_r =
\SI{5}{\electronvolt\per\angstrom\squared}$ and $r_{\mathrm{eq}} =
\SI{1.53}{\angstrom}$, angles with $\kappa_\theta =
\SI{2.5}{\electronvolt\per\radian\squared}$ and $\theta_{\mathrm{eq}} =
\ang{112}$, and a torsion with $\kappa_1 = 0.06$, $\kappa_2 = -0.01$ and
$\kappa_3 = \SI{0.12}{\electronvolt}$, which puts the lowest well at
$\psi = \ang{180}$ (trans), two gauche wells \SI{36.4}{\milli
\electronvolt} higher at $\pm\ang{63.7}$, barriers of
\SI{127.8}{\milli\electronvolt} at $\pm\ang{118.1}$ and
\SI{180}{\milli\electronvolt} at \ang{0}. The constants are chosen for
illustration. The forces of \texttt{mdlab.forcefield.BondedModel} match
finite differences of its energy to \SI{1e-8}{\electronvolt\per\angstrom}.

With no other terms, the density of $\psi$ is exactly $\mathrm{e}^{-\beta
U_{\mathrm t}(\psi)}$, $U_{\mathrm t}$ the torsion energy. Building the
chain bead by bead, the second bead's position relative to the first is
given by a length and two angles, the third's relative to the second by
$r_2$, the bond angle $\theta_1$ and an angle about the first bond, and
the fourth's by $r_3$, $\theta_2$ and its angle about the middle bond,
which is $\psi$. Each step is a set of spherical coordinates: small
changes of $r$, of $\theta$ and of the angle $\phi$ about the axis move the
bead by $\dd r$, $r\,\dd\theta$ and $r\sin\theta\,\dd\phi$ at right angles
to one another, a small volume $r^2\sin\theta\,\dd r\,\dd\theta\,\dd\phi$.
The volume for the whole chain is a product of such factors of $r$ and
$\sin\theta$ times $\dd\psi$, with no factor that depends on $\psi$. The energy is a sum of a part in the bonds and angles
and a part in $\psi$, so integrating over everything but $\psi$ leaves
$\mathrm{e}^{-\beta U_{\mathrm t}(\psi)}$ times a constant, and $F(\psi) =
U_{\mathrm t}(\psi)$ up to a constant.

Two hundred copies, each started in the trans well, are held by Langevin
friction of \SI{1}{\per\pico\second} at \SI{300}{\kelvin} and at
\SI{200}{\kelvin} for \SI{2}{\nano\second}, $\psi$ recorded every
\SI{100}{\femto\second}, and the first \SI{200}{\pico\second} left out:
over it the trans well holds 0.720 and 0.903 of the copies, against the
0.664 and 0.801 of equilibrium. Each bin of $F(\psi)$ is compared with
$-\kB T\ln$ of the bin's average of $\mathrm{e}^{-\beta U_{\mathrm t}}$,
since a histogram counts the whole bin. At \SI{300}{\kelvin}, where each
copy changes well 83.5 times per nanosecond, $F$ matches the torsion
energy within its errors (\cref{fig:ob-fes}c): the mean of the squared
misses over squared errors is 1.02 over 72 bins with blocks in time and
0.68 with blocks of 20 copies each, and at $\psi \approx 0$ the run gives
$(181 \pm 2)$\,\si{\milli\electronvolt} against 179.

At \SI{200}{\kelvin} each copy changes well only 8.0 times per
nanosecond, and the profile is harder to get right. Its shape is right,
$F(0) = (177 \pm 4)$\,\si{\milli\electronvolt} against 179, but with blocks
in time the squared misses over squared errors average 3.26, and they lie
mostly on one side, where counts are plentiful: 7.8 between $-\ang{120}$
and 0, 1.3 between 0 and \ang{120}. Over the \SI{1.8}{\nano\second} kept,
the gauche well near $+\ang{64}$ holds $0.0158 \pm 0.0088$ more of the
copies than the one near $-\ang{64}$, though by symmetry the two must hold
equal shares. That is
1.8 of its errors, an ordinary fluctuation when the error is taken over
the 200 independent copies; but each copy stays in a well for
\SI{125}{\pico\second} on average, so the populations of consecutive blocks in time
are not independent and their spread understates the error, by the mechanism of
\cref{sec:cv-blocks}. Blocks of copies instead give errors 40\% larger
(a median of 0.92 against \SI{0.66}{\milli\electronvolt}) and a mean of
2.05, much of what remains being that single fluctuation, which moves
every bin of a well together: folding $\psi$ to $\lvert\psi\rvert$, which
pools the two gauche wells and so cancels it, brings the mean to 1.06
with blocks of copies, while blocks in time still give 2.10. Free-energy profiles across barriers
crossed rarely need their error bars from independent runs, or from
blocks far longer than the time between crossings.

\begin{practice}
A free-energy profile is no better than the number of crossings between
its wells, and counting them is the check. For a two-well process with independent, approximately exponential stays,
$n$ crossings mean each side
is visited about $n/2$ times, its share of the time known to about
$\sqrt{2/n}$ of itself, so the ratio of the two shares is known to about
$2/\sqrt n$, and the difference of free energy to $2\kB T/\sqrt n$: two
thirds of $\kB T$ for ten crossings and a tenth for 400, whatever the
number of frames. Copies started in different wells, errors taken over
independent copies, and the start left out, where the wells still hold
the populations they began with, guard against the rest.
\end{practice}

\subsection*{The distribution of the potential energy}

The density of the potential energy itself is a profile of the same kind,
and the TrajCast paper compares it between a learned model and the
reference\cite{Thiemann2026}. It depends on the ensemble as well as the
forces: Chapter~15's liquid argon spreads by 0.677 under CSVR and by
\SI{0.532}{\milli\electronvolt} per atom at fixed energy, and the two
distributions overlap by 0.882 by the score of \cref{sec:ob-vdos}, against
0.980 for the two halves of the same CSVR run. A test of a model by this
distribution is a test of its forces only when both runs sample the same
ensemble; with that settled, a score near the halves' own is the target.

\begin{figure}[htb]
\centering
\includegraphics{ch16_observables/fes}
\caption{(a) Lithium on the model surface at \SI{600}{\kelvin}, positions
folded into one cell of the hollows: $F(x, y)$ (shading, \SIrange{0}{0.35}{
\electronvolt}, with its additive constant aligned to $U$) and contours of $U$ every \SI{0.1}{\electronvolt} (lines).
(b) The same atoms' free energy against the distance $r_{\mathrm h}$ to the nearest
hollow (points, with twice their errors), against the exact
$-\kB T\ln[r_{\mathrm h}\oint\mathrm{e}^{-\beta U}\dd\theta]$ (line) and $U$ along the
line from a hollow to a bridge (dashed). (c) The four-bead chain's
dihedral angle at 300 and \SI{200}{\kelvin}: $F(\psi)$ with twice its
error from ten blocks of 20 copies, after the first
\SI{200}{\pico\second}, against the torsion energy (line).}
\label{fig:ob-fes}
\end{figure}

\begin{keyidea}
The free energy along a coordinate is $-\kB T\ln$ of the density of its
values, estimated by a histogram, with errors from blocks. It equals the
potential energy only when the coordinate is every coordinate; otherwise
it counts the arrangements at each value. Across barriers crossed rarely,
short blocks in time underestimate errors; use sufficiently long blocks
or independent copies.
\end{keyidea}

\notebook{16}{build F from a histogram and watch the error bars at 200 K}

\begin{selfcheck}
Why does $F(r_{\mathrm h})$ for the lithium atom rise towards $r_{\mathrm h} = 0$ although the
potential energy is lowest there?
\end{selfcheck}
